\section{Further bounds and limitations}

\begin{proposition}[Unrestricted depth three]\label{prop:c3}
For the infimum $c_3$ of $\Vstar/U$ over laminar forests of depth at
most three with $U>0$,
\[
 \frac{1+4/\pi}{3}\le c_3\le\gamma .
\]
\end{proposition}
\begin{proof}
After taking weight magnitudes the Gram recursion admits a maximizing
system with nonnegative coordinates: place child systems on disjoint
coordinate sets and align each root with the nonnegative sum, using
a fresh coordinate at zero. Let $C$ be this covariance matrix and let
$D=\sum_iw_i^2$. For a Gaussian vector with covariance $C$, each
coordinate is independent of the joint vector of coordinates matured
at its issue because its cross-covariances with this vector vanish.
Gaussian signs therefore satisfy the history LP. The Gaussian sign
correlation formula gives $(2/\pi)\arcsin C_{ij}$. Since $C_{ij}\ge0$,
the resulting objective is at least $D+(2/\pi)(U-D)$.

For a subtree of depth $d$ let $D_v$ be its sum of squared weights.
Inductively $u_v^2\le dD_v$. At a nonleaf use
$\sum_c u_c^2\le(d-1)(D_v-w_v^2)$ and apply
Cauchy--Schwarz to the two terms of the Gram recursion. This gives
$U\le3D$ in depth three, and the lower bound follows. The upper bound
uses the attaining weights of Theorem~\ref{thm:seven}.
\end{proof}

The equality $c_3=\gamma$ is open. The three-star theorem restricts
both the number of child stars and the number of active leaves.
The depth-three formula still follows from the general recursions
with either restriction removed, but its first moment and Gram norm
then need a new joint comparison.

\subsection{Four-star partial results}

The established four-star analysis covers a relative neighbourhood of
the three-star boundary, four identical effective pairs, and the slice
with two zero lighter coefficients. Its explicit local estimate is useful
for identifying an obstruction to a naive deletion argument. Choose two
pairs $(x+u,y+v),(x-u,y-v)$ and let $Z$ be the sum of the remaining
pairs, with Gram norm $B_Z$. Set
\[
 B_0=\sqrt2(x+ky),\quad D_0=2(u+kv)^2+B_Z^2,\quad
 \sigma=\tfrac14\sqrt{c_2/2},\quad \lambda=\sigma/\gamma .
\]
Here $B^2=B_0^2+D_0$. Relative to the two identical pairs $(x,y)$,
the increase in the first moment is at least $\sigma\sqrt{D_0}$.
To see this, condition on the cancellation of the two large signs and
the two small signs. On the joint cancellation event of probability
$1/4$, use sharp Khintchine on $2u\epsilon+2v\eta+Z$.
On the other events Jensen bounds the contribution below by the
corresponding reference contribution. The second moment increases by
at least $c_2D_0$ by Cauchy--Schwarz. Consequently
$\sqrt{D_0}/B_0\le 2\lambda/(1-\lambda^2)$ for $B_0>0$ gives
$A-A_0\ge\gamma(B-B_0)$, and the deficit at $\gamma$ cannot decrease.
On the normalized four-pair domain $B=1$, the minimized deficit is
continuous. On the boundary with a zero full pair it is positive
except at the equality rays of Theorem~\ref{thm:three}, where the
explicit estimate applies. Compactness yields a uniform neighbourhood
of that boundary with an existential radius.

The identical-pair slice has exact minimum
$367/(1919-1040\sqrt2)>\gamma$ in the frozen four-star analysis.
For four copies of $(a,b)$, the magnitude of the sum of four
independent fair signs has masses $3/8,1/2,1/8$ at $0,2,4$.
Using two independent copies gives
$A=3a/2+9b/16+(2b-a)_+/8$, $Q=4(a^2+b^2)$ and $B=2(a+kb)$.
Discarding the positive-part term gives a quadratic quotient
with numerator matrix and denominator vector
\[
 M_4=\begin{pmatrix}1&3/2&9/16\\3/2&4&0\\9/16&0&4\end{pmatrix},
 \qquad \ell_4=(1,2,2k)^\top .
\]
Its Schur complement is $367/1024>0$. The lower bound is
$(\ell_4^\top M_4^{-1}\ell_4)^{-1}=367/(1919-1040\sqrt2)$.
Direct substitution gives positive coordinates for $M_4^{-1}\ell_4$
with $a>2b$, so the discarded term vanishes at the minimizing ray.
Two singleton pairs $(a,0),(b,0)$ may be merged into one sorted pair
without changing the Rademacher law or decreasing the Gram norm.
Theorem~\ref{thm:three} then applies. General deletion is obstructed.
For instance, appending $(1/2,1/4)$ to $(5,2),(4,2),(2,0)$ leaves
$A$ unchanged because the old sum is odd and the added magnitude
is at most $3/4$. At the old minimizing outer weight the increase
in the Gram denominator is sufficient to decrease the quotient.
The companion source map identifies the exact comparisons in S08
and the reproduced fixture. The resulting quotient remains above
$\gamma$.

Selected closed four-star chambers have separate exact certificates.
The sharp child of chamber C illustrates a distinction between the
variance inequality and its certificate. In the fixed-projection
quadratic comparison a zero on a two-ray support imposes a kernel
condition on the PSD part of the aggregate residual. The certified
residual at this split and parameter $\delta=56/25$ requires rank two
in a PSD-plus-entrywise-nonnegative decomposition. The other child has
a certificate of rank one after the projection is changed. Their union
covers C. This is a local result for specified geometric data.

Copositive zero supports, SPN decompositions and simplicial verification
are established machinery \cite{hildebrand,shaked,bundfuss,bundfussadaptive}. A negative
value of one aggregate on a nonnegative vector excludes an SPN
certificate for that aggregate at any rank. The majorization can be
changed while retaining the original variance problem. For the
polynomial comparisons used with shifted projections, the Bernstein
basis is nonnegative and sums to one on the closed interval
\cite{funfzig}. Nonnegative coefficients give a nonnegative polynomial.
Strict positivity of an anchor requires a strictly positive lower
bound on the whole interval. All strictly positive coefficients give
such a bound, as in the two frozen positive examples. For instance,
the polynomial $t$ has nonnegative coefficients on $[0,1]$ and
vanishes at the endpoint. The full four-star inequality remains open.

\subsection{Interpretation and limitations}

The sharp comparisons concern two optima. Every admissible process
satisfies $\Var(\sum_i w_iX_i)\le\Vstar\le U$, and the identically
zero scores show that a positive multiple of $U$ gives no variance
floor for an individual process. A prescribed outcome mechanism or
richer fixed filtration can shrink the attainable class. Adaptive
weights, random schedules, and intervals that cross without nesting
fall outside the stated sharp laminar theorems. Marginal centering
alone cannot replace issue-time centering. These results supply no
forecast accuracy, sequential-testing or profitability conclusion.
